\section{Decoding and Measurement}
\label{sec:decoding}

The single-qubit decoding formula~\eqref{eq:decoding} recovers
the encoded value from the amplitude ratio of the output state.
In the multi-qubit setting, the system register is generically in
a superposition of eigenspaces, and measuring the ancilla affects
the system state.  This section analyzes what $f(H)$ means in the
stereographic framework, characterizes the post-measurement state,
and quantifies the measurement overhead.

\subsection{Eigenspace Superposition}
\label{subsec:superposition}

Suppose the initial system state is a superposition
$\ket{\psi} = \sum_j c_j \ket{j}$, where $\ket{j}$ are
eigenstates of~$H$ with eigenvalues~$\lambda_j$.  After the
stereographic QSVT protocol
(Theorem~\ref{thm:eigenvalue-transform}), the joint state of
ancilla and system is
\begin{equation}
	\mathcal{W}_{\bm{\phi}}\ket{0}_a\ket{\psi}_s
	= \sum_j c_j
	\bigl(P(\tilde{\lambda}_j)\ket{0}
	+ Q(\tilde{\lambda}_j)\ket{1}\bigr)\ket{j}_s.
	\label{eq:superposition-state}
\end{equation}
This is an entangled state of ancilla and system: the ancilla
state depends on which eigenspace the system occupies.

\subsection{Post-Measurement State}
\label{subsec:post-measurement}

Measuring the ancilla in the computational basis collapses the
system into a conditional state.  Projecting onto
$\ket{0}_a$ gives the (unnormalized) state
\begin{equation}
	\ket{\psi_0} = \sum_j c_j\, P(\tilde{\lambda}_j)\,\ket{j}_s,
	\label{eq:post-measurement-0}
\end{equation}
with probability
$p_0 = \sum_j |c_j|^2\,|P(\tilde{\lambda}_j)|^2$.
Projecting onto $\ket{1}_a$ gives
\begin{equation}
	\ket{\psi_1} = \sum_j c_j\, Q(\tilde{\lambda}_j)\,\ket{j}_s,
	\label{eq:post-measurement-1}
\end{equation}
with probability $p_1 = 1 - p_0$.

\begin{proposition}[Post-Measurement Interpretation]
	\label{prop:post-measurement}
	The post-measurement state $\ket{\psi_0}$ (conditioned on
	ancilla outcome $\ket{0}$) is the result of applying the
	operator $P(\tilde{H})$ to $\ket{\psi}$, where
	$\tilde{H} = H/\sqrt{\Id + H^2}$ is the
	stereographically compressed Hamiltonian (defined via
	functional calculus).  Similarly, $\ket{\psi_1}$
	corresponds to $Q(\tilde{H})\ket{\psi}$.
\end{proposition}

This means stereographic QSVT implements a \emph{probabilistic}
polynomial transformation of~$H$: with probability $p_0$, the
system is transformed by $P(\tilde{H})$, and with probability
$p_1$, by $Q(\tilde{H})$.

\subsection{Extracting the Rational Function}
\label{subsec:extracting}

The decoded rational function $f(\lambda_j) =
	P(\tilde{\lambda}_j)/Q(\tilde{\lambda}_j)$ is extracted from
measurements of the ancilla Pauli operators.  For an input
eigenstate $\ket{j}$, the decoding formula~\eqref{eq:decoding}
applies directly.  For a superposition input, one obtains the
\emph{expectation value}
\begin{equation}
	\frac{\langle\PauliX\rangle_a}
	{1 - \langle\PauliZ\rangle_a}
	= \frac{\sum_j |c_j|^2\,\re[P\bar{Q}]\,\sqrt{1+\lambda_j^2}}
	{\sum_j |c_j|^2\,|Q|^2},
	\label{eq:superposition-decoding}
\end{equation}
which is a weighted average over eigenspaces, \emph{not}
$f(\lambda_j)$ for any single eigenvalue.  To extract
$f(\lambda_j)$ for individual eigenvalues, one must either:
\begin{enumerate}
	\item prepare the input in a single eigenstate $\ket{j}$
	      (which requires eigenbasis knowledge), or
	\item combine the protocol with eigenvalue-resolving
	      techniques such as quantum phase estimation.
\end{enumerate}

\subsection{Measurement Overhead}
\label{subsec:overhead}

The measurement overhead for stereographic decoding was analyzed
in the companion paper~\cite{KharaziLeytonOrtega2026a}.  For a
single eigenvalue $r = \lambda_j$, achieving absolute precision
$\epsilon$ in the decoded value requires
\begin{equation}
	N = \mathcal{O}\!\left(\frac{(1+r^2)^2(1+r)^2}{\epsilon^2}\right)
	= \mathcal{O}\!\left(\frac{r^6}{\epsilon^2}\right)
	\quad \text{for } r \gg 1
	\label{eq:shot-count}
\end{equation}
measurement shots.  For relative precision the scaling improves
to $\mathcal{O}(r^4/\epsilon_{\mathrm{rel}}^2)$.

In the multi-qubit setting, the relevant quantity is the
\emph{maximum} eigenvalue being decoded:
$r_{\max} = \max_j \lambda_j$.  The total measurement budget
scales as $\mathcal{O}(r_{\max}^6/\epsilon^2)$ for absolute
precision, compared to $\mathcal{O}(1/\epsilon^2)$ for standard
QSVT.

% TODO: Add analysis of adaptive measurement strategies
% that could improve the scaling for specific eigenvalue ranges
